% ======================================================================
% Chapter 3: Scope of Lattice QCD Calculations
% Source: R. Gupta, "Introduction to Lattice QCD" (arXiv:hep-lat/9807028),
% Section 3, pp. 10-14. The closing paragraph and bullet list continue onto
% page 14 (their text is extracted at the top of sec04.txt).
% ======================================================================

\chapter{Scope of Lattice QCD Calculations}

\begin{figure}[htbp]
\centering
\includegraphics[width=0.8\textwidth]{images/fig1.png}
\caption{The CKM matrix in the Wolfenstein parameterization. I show examples
of physical processes that need to be measured experimentally and compared
with theoretical predictions via the ``master'' equation to estimate various
elements of the CKM matrix.}
\label{fig:1}
\end{figure}

The primary goal of LQCD is to test whether QCD is the correct theory of
strong interactions. Most particle physicists believe that the answer is
yes because of its successes in describing processes with large momentum
transfer, its mathematical elegance, and uniqueness. However, quantitative
confirmation is lacking. The problem with the conventional perturbative
approach (expansion in a small parameter) when analyzing hadronic process
at scales $\lesssim 1$~GeV is that the strong coupling constant
$\alpha_s \sim 1$. Thus, perturbation theory in $\alpha_s$ is not reliable.
As a result we cannot calculate the masses of mesons and baryons from QCD
even if we are given $\alpha_s$ and the masses of quarks. This is
illustrated in Fig.~\ref{fig:2} where I show a comparison between the
binding energy of a hydrogen atom and that of the proton. It turns out that
almost all the mass of the proton is attributed to the strong non-linear
interactions of the gluons.

\begin{figure}[htbp]
\centering
\includegraphics[width=0.8\textwidth]{images/fig2.png}
\caption{An illustration of the difference in the binding energy between the
electron and the proton in a hydrogen atom (interacting via electromagnetic
forces) and the quarks in a proton (strong force).}
\label{fig:2}
\end{figure}

Simulations of lattice QCD have six unknown input parameters. These are the
coupling constant $\alpha_s$ and the masses of the up, down, strange, charm
and bottom quarks (the top quark is too short lived, $0.4 \times 10^{-24}$
seconds, to form bound states and is thus not suitable for lattice
studies). Once these have been fixed in terms of six precisely measured
masses of hadrons, thereafter, the masses and properties of the hundreds of
other particles made up of these quarks and gluons have to agree with
experiments. For example, consider the chiral perturbation theory
($\chi$PT) relation for the pion mass
\begin{equation}
(M_{\pi^+})^2 = B_\pi (m_u + m_d) + C_\pi (m_u + m_d)^2 + \ldots
\label{eq:chipt}
\end{equation}
Lattice calculations will provide a check of this relation and fix the
constants $B_\pi$, $C_\pi$, \ldots, i.e.\ constrain the chiral lagrangian.
In addition, having measured $M_\pi$ for a number of values of the quark
mass, one can invert Eq.~\eqref{eq:chipt} to determine the current quark
masses. I shall call such tests of QCD and the determination of the six
unknown parameters of the SM a direct application of lattice QCD.

The second area in which perturbative estimates are not reliable is in the
calculation of the matrix elements occurring in the weak decays of hadrons.
The reason is that the non-perturbative QCD corrections to the basic weak
process can be large due to the exchange of soft gluons between the initial
and final states. This is illustrated in Fig.~\ref{fig:3} for the case of
the semi-leptonic decay $D \to K l\nu$. In this case the matrix element of
the weak interactions Hamiltonian between the initial $D$ meson and final
kaon receives large corrections from QCD which cannot be estimated using
PQCD.

\begin{figure}[htbp]
\centering
\includegraphics[width=0.8\textwidth]{images/fig3.png}
\caption{The Feynman diagram for the semi-leptonic decay
$D^- \to K^0 e^- \bar{\nu}_e$. The QCD corrections are illustrated by the
various gluons being exchanged between the initial and final hadrons. The
leptonic vertex can be calculated reliably using perturbation theory,
whereas the hadronic vertex requires non-perturbative methods.}
\label{fig:3}
\end{figure}

The result of the experiment is a number, i.e.\ the decay rate. To check
whether the Standard Model prediction is in agreement with this number
theorists must derive an analytical expression for the same process. This
expression, in general, consists of a product of three parts: a combination
of the parameters in the SM, matrix elements (ME) of the interaction
Hamiltonian between the initial and final states, and kinematic factors.
(This is the well-known Fermi's golden rule.) Schematically one can write
this ``master'' equation as follows.
\begin{equation}
\mathrm{Expt.}\# =
(\mathrm{SM\ parameters})\,(\mathrm{matrix\ elements})\,(\mathrm{kinematic\ factors}).
\label{eq:master}
\end{equation}
Thus, for each such experimental number one gets a constraint on one
particular combination of the SM parameters provided the ME are known.

By using several precision experiments and calculating the necessary ME one
can over-constrain the model and therefore test it. I shall call lattice
calculations of ME an ``indirect'' probe of Standard Model parameters as
one usually needs additional phenomenological/experimental input to extract
the unknown parameters via the master equation.

Physicists have been busy calculating and measuring these parameters for
the last 25 years, so one already has reasonable estimates. It is,
therefore, useful to keep in mind the level of precision that lattice
calculations have to aim for to make an impact. Non-lattice estimates,
taken from the 1996 Review of Particle Properties and from Prof.~Buras's
lectures [10], are

\begin{table}[htbp]
\centering
\caption*{Non-lattice estimates of the Standard Model parameters}
\label{tab:nl-estimates}
\begin{tabular}{l l p{6.5cm}}
\toprule
Parameters & Value & Comments \\
\midrule
$m_u(\overline{MS}, 2\ \mathrm{GeV})$ & 3--4.5~MeV & $\chi$PT, sum-rules \\
$m_d(\overline{MS}, 2\ \mathrm{GeV})$ & 5--8~MeV & $\chi$PT, sum-rules \\
$m_s(\overline{MS}, 2\ \mathrm{GeV})$ & 100--140~MeV & $\chi$PT, sum-rules \\
$m_c(\overline{MS}, m_c)$ & 1.0--1.6~GeV & $J/\psi$ sprectra \\
$m_b(\overline{MS}, m_b)$ & 4.1--4.6~GeV & $\Upsilon$ spectra \\
$M_{\mathrm{Higgs}}$ & & Not Found \\
$\alpha_s(M_Z)$ & $0.118 \pm 0.003$ & World average \\
$|V_{ud}|$ & $0.9736 \pm 0.0010$ & \begin{tabular}[t]{@{}l@{}}
$n \to p e\nu_e$ and \\ $\mu \to e\nu_e\nu_\mu$ \end{tabular} \\
$|V_{us}|$ & $0.2196 \pm 0.0023$ & \begin{tabular}[t]{@{}l@{}}
$K^+ \to \pi^0 e^+ \nu_e$ and \\ $K_L^0 \to \pi^\pm e^\mp \nu_e$ \end{tabular} \\
$|V_{ub}|$ & $(3.3 \pm 0.2 \pm 0.4 \pm 0.7) \times 10^{-3}$ &
\begin{tabular}[t]{@{}l@{}} $B \to \pi l\nu$; $B \to \rho l\nu$ \end{tabular} \\
$|V_{cb}|$ & $0.040 \pm 0.003$ & \begin{tabular}[t]{@{}l@{}}
$B \to X_c l\nu$ and \\ $B \to D^* l\nu$ \end{tabular} \\
$|\epsilon|$ & $2.26 \times 10^{-3}$ & $K^0 \leftrightarrow \bar{K}^0$ mixing \\
$|\epsilon/\epsilon'|$ & $0$--$30 \times 10^{-4}$ & $K \to \pi\pi$ decays \\
\bottomrule
\end{tabular}
\end{table}

Note that only the magnitude of the CKM elements can be determined from
experiments. Clearly, the most enticing opportunity for lattice QCD is to
help determine the quark masses, $\alpha_s$, and the essentially unknown
Wolfenstein parameters $\rho$ and $\eta$. For other elements of $V_{CKM}$
one has reasonable estimates, however here too further progress requires
significant improvement in the estimates of the hadronic matrix elements.
Phenomenologists are relying on LQCD to provide improved estimates for
these.

Let me end this section by mentioning that in addition to validating QCD as
the correct theory of strong interactions by reproducing the measured
spectrum of mesons, baryons, and glueballs, and the calculation of weak
matrix elements, simulations of LQCD will also allow us to
\begin{itemize}
\item Investigate the topological structure of the QCD vacuum, and the
mechanisms responsible for confinement and spontaneously broken chiral
symmetry. For a recent review see [12].
\item Calculate hadronic properties like wavefunctions, decay constants,
form factors for electromagnetic, weak, and rare decays, and structure
functions for deep inelastic scattering. Estimate the pion nucleon sigma
term and the proton spin carried by the quarks by calculating the matrix
elements $\langle N | q\bar{q} | N \rangle$ and
$\langle N | q\gamma_\mu \gamma_5 q | N \rangle$. Some of these topics
will be covered by Prof.~Martinelli, for others a good reference is the
proceedings of LATTICE conferences [9].
\item Analyze the behavior of QCD at high temperatures and address questions
like the nature of the phase transition, the transition temperature $T_c$,
the equation of state, and the collective excitations near and above $T_c$.
The topic of QCD at finite temperature has been covered by J-P Blaizot at
this school, so I will only list for you my favorite pedagogical reviews of
the subject [13]. The lectures by Profs.~Martinelli, L\"{u}scher, and me
will be confined to understanding and probing the properties of QCD at zero
temperature.
\end{itemize}
